\documentclass{article}
\usepackage[UTF8]{ctex}
\usepackage{amsmath,amssymb,amsfonts}
\usepackage{enumitem}
\usepackage{geometry}
\geometry{a4paper, margin=1in}
% === 错题本排版优化（不参与题目标识） ===
\setlength{\parindent}{0pt}
\setlength{\parskip}{0.35em plus 0.10em minus 0.08em}
\setlength{\emergencystretch}{3em}
\setlist{itemsep=0.25em, topsep=0.40em, parsep=0pt}
\allowdisplaybreaks[2]
% === 错题本排版优化结束 ===
\begin{document}
\section*{东北大学2024年高等代数真题}
\section*{65．东北大学2024年高等代数试卷}

\subsection{题目1}
1．（15 分）设有两个向量组 （I）：\(\alpha_{1}=(1,0,2)^{\mathrm{T}}, \alpha_{2}=(1,1,3)^{\mathrm{T}}, \alpha_{3}=(1,-1, a+2)^{\mathrm{T}}\) ；（II）：\(\beta_{1}=(1,2, a+3)^{\mathrm{T}}, \beta_{2}=\) \((2,1, a+6)^{\mathrm{T}}, \beta_{3}=(2,2, a+2)^{\mathrm{T}}\) ．\\
讨论当 \(a\) 为何值时，向量（I）和（II）等价？当 \(a\) 为何值时，向量组（I）能由向量组（II）线性表出，但向量组（II）不能由向量组（I）线性表出？

\subsection{题目2}
2．（15 分）计算 \(n\) 阶行列式

\[
D_{n}=\left|\begin{array}{ccccc}
x & a & a & \cdots & a \\
b & x & a & \cdots & a \\
b & b & x & \cdots & a \\
\vdots & \vdots & \vdots & & \vdots \\
b & b & b & \cdots & x
\end{array}\right| .
\]

\subsection{题目3}
3．（15 分）设实二次型 \(f\left(x_{1}, x_{2}, x_{3}\right)=\left(x_{1}+\lambda_{1} x_{2}\right)^{2}+\left(x_{2}+\lambda_{2} x_{3}\right)^{2}+\left(x_{3}+\lambda_{3} x_{1}\right)^{2}\) ．

\begin{itemize}
  \item[（1）] （5 分）当 \(\lambda_{1}, \lambda_{2}, \lambda_{3}\) 满足什么条件时，\(f\left(x_{1}, x_{2}, x_{3}\right)\) 是正定二次型？
  \item[（2）] （10 分）若 \(f\left(x_{1}, x_{2}, x_{3}\right)\) 经过正交线性替换可化为标准形 \(y_{1}^{2}+y_{2}^{2}+y_{3}^{2}\) ，求 \(\lambda_{1}, \lambda_{2}, \lambda_{3}\) 的值．
\end{itemize}

\subsection{题目4}
4．（15 分）设 \(f(x)=a+b x\) 是有理数域上的一元多项式，\(\alpha=(1,1, \cdots, 1)^{\mathrm{T}}\) 和 \(\beta=(-1,-1, \cdots,-1)^{\mathrm{T}}\)是 \(n\) 维列向量。令 \(J=\beta \alpha^{T}, A=f(J)\) 。

\begin{itemize}
  \item[（1）] （5 分）求 \(J\) 的全部特征值和属于每个特征值的全部特征向量．
  \item[（2）] （10 分）\(A\) 是否能相似于一个对角矩阵？如果能，求一个可逆矩阵 \(P\) ，使得 \(P^{-1} A P\) 为对角矩阵，并写出这个对角矩阵．
\end{itemize}

\subsection{题目5}
5．（15 分）已知

\[
A=\left(\begin{array}{ccc}
0 & 1 & 0 \\
0 & 0 & 1 \\
-2 & 3 & -1
\end{array}\right) .
\]

分别在有理数域和复数域上讨论矩阵 \(A\) 是否可对角化？并说明理由。

\subsection{题目6}
6．（15 分）设 \(V\) 是 \(\mathbb{R}\) 上所有二阶矩阵构成的线性空间，\(V\) 上的两组基为

\[
\begin{aligned}
& E_{11}=\left(\begin{array}{ll}
1 & 0 \\
0 & 0
\end{array}\right), E_{12}=\left(\begin{array}{ll}
0 & 1 \\
0 & 0
\end{array}\right), E_{21}=\left(\begin{array}{ll}
0 & 0 \\
1 & 0
\end{array}\right), E_{22}=\left(\begin{array}{ll}
0 & 0 \\
0 & 1
\end{array}\right) ; \\
& B_{11}=\left(\begin{array}{ll}
1 & 0 \\
0 & 0
\end{array}\right), B_{12}=\left(\begin{array}{ll}
1 & 1 \\
0 & 0
\end{array}\right), B_{21}=\left(\begin{array}{ll}
0 & 0 \\
1 & 0
\end{array}\right), B_{22}=\left(\begin{array}{ll}
1 & 1 \\
0 & 1
\end{array}\right) .
\end{aligned}
\]

定义 \(V\) 上的线性变换 \(\mathcal{A}\) 为

\[
\mathcal{A}(X)=\left(\begin{array}{cc}
1 & 3 \\
-2 & -6
\end{array}\right) X, X \in V .
\]

（1）（5 分）求基 \(E_{11}, E_{12}, E_{21}, E_{22}\) 到 \(B_{11}, B_{12}, B_{21}, B_{22}\) 的过渡矩阵。\\
（2）（10 分）求线性变换 \(\mathcal{A}\) 在基 \(B_{11}, B_{12}, B_{21}, B_{22}\) 下的矩阵。

\subsection{题目7}
7．（15 分）设 \(f(x)=\left(x-k_{1}\right)\left(x-k_{2}\right) \cdots\left(x-k_{n}\right)+1\) ，其中 \(k_{1}, k_{2}, \cdots, k_{n}(n>2)\) 是互异的整数。证明：\\
（1）（5 分）多项式 \(f(x)\) 在实数域上可约．\\
（2）（10 分）若多项式 \(f(x)\) 在有理数域上可约，则 \(n\) 必为偶数，且 \(f(x)\) 必能表示成某个整系数多项式的平方．

\subsection{题目8}
8．（15 分）设 \(A\) 为 \(n\) 阶正定矩阵，证明：存在下三角矩阵 \(B\) ，使得 \(A=B B^{\mathrm{T}}\) 。

\subsection{题目9}
9．（15 分）证明如下问题：

\begin{itemize}
  \item[（1）] （5 分）\(n \times n\) 数字矩阵 \(A\) 的特征矩阵是满秩矩阵，但不是可逆矩阵。
  \item[（2）] （10 分）\(n \times n\) 数字矩阵 \(A\) 是可逆矩阵的充要条件是 \(A\) 的最后一个不变因子 \(d_{n}(\lambda)\) 有非零常数项．
\end{itemize}

\subsection{题目10}
10．（15 分）设 \(A\) 为 \(n\) 阶方阵，且 \(A^{3}-5 A^{2}+4 A=O\) 。证明：

\[
\operatorname{rank}(A)=\operatorname{rank}\left(A^{2}-A\right)+\operatorname{rank}\left(A^{2}-4 A\right)
\]
\end{document}
